\subsection{Antecedentes: conceptos de base}
\label{anx:antecedentes}

Esta sección desarrolla los conceptos de la Tabla~\ref{tab:antecedentes}.

\subsubsection{Sistemas de recomendación}

Un sistema de recomendación es un software que sugiere a una persona los elementos que con
mayor probabilidad le resultan útiles, a partir de información sobre sus necesidades, sus
preferencias o el comportamiento de otros usuarios \citep{adomavicius_toward_2005}. La
literatura distingue varios enfoques. Los sistemas basados en contenido recomiendan
elementos parecidos a los que la persona ya eligió, los de filtrado colaborativo recomiendan
lo que eligieron usuarios con gustos similares, y los basados en conocimiento razonan sobre
las necesidades explícitas del usuario y sobre las características de cada producto
\citep{burke_hybrid_2002}. Los sistemas híbridos combinan varios de estos enfoques para
aprovechar sus fortalezas.\par

En la compra escolar, el filtrado colaborativo resulta poco útil, porque una familia compra
útiles una vez al año y su lista depende del colegio y del grado, no de lo que compraron
otras familias. Por esta razón, Tukit se apoya en el enfoque basado en conocimiento, ya que
recomienda a partir de los requerimientos de la lista oficial y de las preferencias que
declara el acudiente. A diferencia de la mayoría de los sistemas, que recomiendan productos
individuales, Tukit recomienda conjuntos completos, un problema conocido como recomendación
de paquetes, en el que importa la coherencia entre los elementos que se ofrecen juntos
\citep{avny_brosh_bruce_2022}. Como cada paquete debe equilibrar varios criterios al mismo
tiempo, el sistema se clasifica como un sistema de recomendación multiobjetivo.\par

\subsubsection{Comprensión de documentos y modelos multimodales}

El primer reto, leer la lista, pertenece al campo de la comprensión de documentos, que
estudia cómo convertir una imagen o un archivo en información estructurada. El enfoque
tradicional es el reconocimiento óptico de caracteres (OCR), que identifica en la imagen las
zonas con texto y reconoce cada carácter \citep{smith_overview_2007}. Este enfoque funciona
bien con documentos impresos y nítidos, pero lee línea por línea sin comprender la
estructura de la página, de modo que no distingue si dos líneas pertenecen a la misma fila
de una tabla o si una viñeta continúa el renglón anterior.\par

Los modelos multimodales de lenguaje superan esa limitación, porque reciben al mismo tiempo
una imagen y una instrucción escrita y razonan sobre la disposición y el contenido de la
página en conjunto. Evaluaciones recientes muestran que, a partir de la imagen únicamente,
estos modelos alcanzan un desempeño comparable al de los flujos que combinan OCR y modelos
de lenguaje \citep{shen_ocr_2026}, aunque su precisión disminuye con el texto manuscrito
\citep{liu_ocrbench_2024}. Para reducir el riesgo de que el modelo produzca información que
no está en el documento, se le exige responder con un formato fijo que se verifica antes de
aceptarse. Este es el enfoque que utiliza Tukit para leer listas en fotografía, PDF o
texto.\par

\subsubsection{Representación vectorial y búsqueda híbrida}

El segundo y el tercer reto, entender cada renglón y encontrar su producto, corresponden a un
problema de búsqueda. La búsqueda léxica compara las palabras de dos textos y tolera
errores de ortografía mediante la distancia de edición, que cuenta cuántas letras deben
cambiarse para pasar de una palabra a otra. Su limitación es que no reconoce sinónimos, de
modo que ``estuche'' y ``cartuchera'' le resultan completamente distintos.\par

La búsqueda semántica resuelve ese problema mediante la representación vectorial del texto.
Un modelo de lenguaje convierte cada frase en un vector, conocido como \emph{embedding}, que
funciona como un conjunto de coordenadas en un espacio donde las frases con significado
parecido quedan cerca \citep{reimers_sentence-bert_2019}. La cercanía entre dos vectores se
mide con la similitud coseno, y a este tipo de búsqueda se le denomina recuperación densa
\citep{karpukhin_dense_2020}. Sin embargo, la búsqueda semántica presenta el riesgo
contrario al de la léxica, ya que puede considerar parecidos dos textos que tratan del mismo
tema sin pedir lo mismo.\par

Estas técnicas suelen integrarse en la generación aumentada por recuperación (RAG), en la
que un modelo de lenguaje redacta una respuesta a partir de los documentos recuperados
\citep{lewis_retrieval-augmented_2020}. En Tukit no es necesario redactar ninguna respuesta,
porque el resultado de la búsqueda es el concepto del catálogo y sus productos. Por ello, el
sistema toma de este esquema la indexación y la recuperación, y combina ambas búsquedas de
forma que la léxica determina qué conceptos son admisibles y la semántica desempata entre
ellos o reconoce sinónimos que no están registrados en el catálogo.\par

\subsubsection{Optimización multiobjetivo y algoritmo NSGA-II}

El cuarto reto consiste en elegir un producto para cada renglón de la lista. Si cada renglón
tiene tres productos posibles, una lista de \(n\) renglones admite \(3^{n}\) combinaciones,
una cantidad que crece tan rápido que una lista de cuarenta renglones supera los
\(10^{19}\) kits. Además, no existe un único kit que sea el mejor, porque el más barato no
es el de mejores marcas, y el que incluye los personajes que pidió el niño no es el más
económico. Los problemas en los que varios criterios compiten entre sí se estudian en la
optimización multiobjetivo \citep{ehrgott_multiobjective_2002}.\par

En este tipo de problemas se dice que una solución domina a otra cuando no es peor en ningún
criterio y es mejor en al menos uno. Las soluciones que ninguna otra domina forman el frente
de Pareto, es decir, el conjunto de alternativas en las que solo es posible mejorar un
criterio cediendo en otro. Cuando las combinaciones son pocas, el frente se puede obtener
revisándolas todas; cuando son muchas, se recurre a algoritmos que lo aproximan en un tiempo
razonable. El más utilizado para problemas con pocos objetivos es NSGA-II
\citep{deb_fast_2002}, un algoritmo genético que mantiene una población de soluciones, las
combina y las modifica en cada generación, y conserva las que pertenecen a los mejores
frentes y las que mantienen mayor variedad entre sí. La calidad del frente obtenido puede
medirse con el hipervolumen, que calcula el espacio de criterios que el frente domina
\citep{zitzler_multiobjective_1999}. Tukit utiliza este enfoque para construir sus kits y
toma del frente la alternativa que mejor responde a cada perfil.\par

\subsubsection{Revisión humana en sistemas de inteligencia artificial}

Ninguna de las técnicas anteriores es infalible, y en la compra escolar un error tiene un
costo directo para la familia. Por esta razón, los sistemas de inteligencia artificial que
toman decisiones con consecuencias para las personas suelen incorporar la revisión humana
en los casos de menor certeza, un esquema conocido como \emph{human-in-the-loop}. Las pautas
de interacción entre personas e inteligencia artificial recomiendan que el sistema indique
con claridad lo que no sabe y que facilite la corrección de sus errores
\citep{amershi_guidelines_2019}. En Tukit, los renglones que el sistema no resuelve con
certeza suficiente pasan a una cola de revisión del equipo antes de que la lista llegue a
las familias.\par

\subsubsection{Carga cognitiva y arquitectura de la decisión}

El último reto se relaciona con la forma en que las personas toman decisiones. La teoría de
la carga cognitiva establece que la memoria de trabajo tiene una capacidad limitada, y que
cuando una tarea exige procesar demasiada información a la vez, la calidad de las decisiones
disminuye \citep{sweller_cognitive_2011}. En la misma línea, se ha demostrado que un número
excesivo de opciones puede desalentar la elección en lugar de facilitarla
\citep{iyengar_when_2000}. La arquitectura de la decisión propone que la forma de presentar
las alternativas, es decir, cuántas son, en qué orden aparecen y cuál se sugiere, influye en
la elección sin limitar la libertad de la persona \citep{thaler_nudge_2008}. A partir de
estos principios, Tukit presenta pocos kits bien diferenciados, cada uno con una razón clara
para elegirlo, y permite que el acudiente modifique cualquier producto antes de pagar.\par
